\documentclass[11pt,a4paper]{article}

\usepackage[margin=1in]{geometry}
\usepackage{amsmath,amssymb,amsfonts}
\usepackage{graphicx}
\usepackage{booktabs}
\usepackage{hyperref}
\usepackage{microtype}
\usepackage{cite}
\usepackage{xcolor}

\hypersetup{
    colorlinks=true,
    linkcolor=blue!70!black,
    citecolor=blue!70!black,
    urlcolor=blue!70!black
}

\title{\textbf{DWMA: Developmental World-Model Agent}\\
\large An Embodied Predictive-Coding and Active-Inference Architecture Without an LLM Cognitive Core}

\author{
  \textbf{Kunal Lubhana}\thanks{Corresponding author: \texttt{research@dwma-project.org}} \\
  \textit{Independent Research / Antigravity Cognitive Systems}
}

\date{September 2026 --- Version 1.0 Preprint}

\begin{document}

\maketitle

\begin{abstract}
We present the \textbf{Developmental World-Model Agent (DWMA)}, an embodied computational architecture designed to investigate whether an autonomous agent can acquire, maintain, and update predictive structure about an interactive dynamical environment without a Large Language Model (LLM) as its cognitive core. DWMA combines predictive forward modeling, epistemic active inference, online structural adaptation, sensorimotor concept formation, symbol grounding, and a functional self-model.

The evaluation is strictly bounded to a continuous 2D simulated dynamical environment where the agent has zero privileged access to latent physical constants (mass, friction, or gravity), requiring these to be inferred purely through sensorimotor interaction residuals. We report the complete, audited four-phase empirical evaluation program conducted across $N = 50$ paired deterministic seeds with $B = 10,000$ percentile bootstrap resamples:
\begin{enumerate}
    \item \textbf{Frozen Baseline Suite (BM-1 to BM-7):} Successful replication across 50 seeds ($49/50$ zero-shot transfer, $50/50$ blind navigation, $50/50$ drift recalibration, $50/50$ epistemic inquiry, $50/50$ counterfactual repair, $50/50$ non-privileged composition, and $48/50$ actuator polarity inversion), accompanied by an $8.8\times$ forecasting MSE reduction against a standardized persistence observer ($0.3276$ vs $2.8803$ MSE, Cohen's $d = 11.45$).
    \item \textbf{Cross-World Transfer ($A \to B \to C$):} Clear zero-shot dynamical shock in novel worlds B ($1.3079 \pm 0.0421$) and C ($0.8117 \pm 0.0272$), followed by rapid online structural adaptation in B reaching $E(t) < 0.20$ within $T_\epsilon = 16.32 \pm 0.62$ steps ($50/50$ pass) with consistent adaptation rate $A_{\text{adapt}} = 0.0461 \pm 0.0084$.
    \item \textbf{Continual Retention ($A \to B \to A$):} Multi-schema memory preserves familiar World A dynamics with near-zero degradation ($E_{A1} = 0.00500 \to E_{A2} = 0.00501$, retention index $\mathcal{R} = 0.9901 \pm 0.0141$, $50/50$ pass $\mathcal{R} \ge 0.75$), whereas an overwriting baseline suffers catastrophic forgetting ($E_{A2} = 0.37460$, $\mathcal{R} = -72.97$, $0/50$ pass).
    \item \textbf{Nonlinear Hypothesis Discrimination:} Primary survival analysis demonstrates that epistemic action selection discriminates between linear viscous drag ($H_1$) and quadratic drag ($H_2$) substantially faster and more reliably than undirected exploration (Kaplan-Meier median latency of $7.0$ steps with $0\%$ censoring vs $23.0$ steps with $28.0\%$ right-censored at the 30-step observation horizon; Mantel-Cox log-rank $\chi^2 = 92.28, p < 10^{-20}$; $50/50$ uncensored DWMA runs vs $36/50$ uncensored random runs).
\end{enumerate}
All empirical claims are strictly bounded to the simulated dynamical environment. No claims of AGI, consciousness, metaphysical causality, or human-like understanding are made.
\end{abstract}

\section{Introduction}
An alternative to making language modeling the cognitive core of an autonomous system is to study learning through direct sensorimotor interaction. Biological development motivates the question: useful internal structure can emerge through repeated action, observation, prediction, and correction rather than through explicit access to latent physical constants.

DWMA is a computational investigation of this principle. It does not claim to reproduce biological development. Instead, it asks whether a non-privileged embodied agent can construct predictive and reusable sensorimotor structure inside a controlled dynamical environment.

\section{Architecture \& Mathematical Formulation}
At discrete time step $t$, the ego sensory observation vector is:
\begin{equation}
\mathbf{o}_t^{ego} = [x_t, y_t, v_{x,t}, v_{y,t}, \theta_t] \in \mathbb{R}^5
\end{equation}
Crucially, physical constants are excluded from observation: $m \notin \mathcal{O}, e \notin \mathcal{O}, \mu \notin \mathcal{O}$.
The continuous motor actuation vector is $\mathbf{a}_t = [u_x, u_y] \in [-1.0, 1.0]^2$, yielding physical thrust:
\begin{equation}
\mathbf{F}_t = \alpha \mathbf{W}_{\text{actuator}} \mathbf{a}_t
\end{equation}
where $\alpha = 0.85$ is the base actuator scaling factor.

The predictive world model computes next-state expectations:
\begin{equation}
\hat{\mathbf{s}}_{t+1} = \mathcal{M}(\mathbf{s}_t, \mathbf{a}_t)
\end{equation}
Violations between prediction and observation generate prediction error residuals $E(t) = \|\mathbf{o}_t - \hat{\mathbf{o}}_t\|^2$, driving online parameter updates and triggering structural adaptation when directional alignment accumulates negative evidence:
\begin{equation}
\Lambda_t = \gamma \Lambda_{t-1} + (1 - \gamma) \cos(\theta_t)
\end{equation}

\begin{figure}[htbp]
\centering
\includegraphics[width=0.85\linewidth]{figures/fig1_system_architecture.png}
\caption{\textbf{DWMA Closed-Loop Architecture.} Continuous motor actuation, reafference-cancelled sensory observation, online forward modeling, epistemic active inference, and multi-schema continual memory.}
\label{fig:arch}
\end{figure}

\section{Cross-World Transfer \& Structural Adaptation ($A \to B \to C$)}
To evaluate whether DWMA acquires reusable physical structure rather than local over-fitting, the agent was evaluated across three distinct dynamical worlds:
\begin{itemize}
    \item \textbf{World A (Baseline):} Friction $\mu = 0.88$, base actuator gain $\alpha = 0.85$, polarity $\mathbf{W}_{\text{act}} = +\mathbf{I}$, gravity $g = 9.8$.
    \item \textbf{World B (Alien Dynamics):} Friction $\mu = 0.71$, base actuator gain $\alpha = 0.85$, inverted actuator polarity $\mathbf{W}_{\text{act}} = -\mathbf{I}$, gravity $g = 4.2$.
    \item \textbf{World C (Radical Dynamic Shift):} Friction $\mu = 0.12$, relative actuator matrix multiplier $\mathbf{W}_{\text{act}} = 0.60\mathbf{I}$ (effective simulator net gain $\alpha \mathbf{W}_{\text{act}} = 0.85 \times 0.60 = 0.51$), polarity $= +\mathbf{I}$, gravity $g = 15.7$.
\end{itemize}

\begin{figure}[htbp]
\centering
\includegraphics[width=0.80\linewidth]{figures/fig2_cross_world_adaptation.png}
\caption{\textbf{Phase 2 Online Structural Adaptation in World B.} Empirical error trajectory $E(t)$ showing transfer shock, autonomous polarity inversion at step 7.0, threshold crossing at $T_\epsilon = 16.32 \pm 0.62$ steps, and asymptotic re-settling to $E(T) = 0.0040$.}
\label{fig:adaptation}
\end{figure}

\section{Continual Retention \& Catastrophic Forgetting ($A \to B \to A$)}
The continual retention benchmark tests whether intermediate adaptation to World B induces catastrophic destruction of World A latent representations. Full DWMA utilizes a multi-schema memory structure that isolates distinct dynamical contexts into discrete schemas $\mathcal{S}_k$.

\begin{figure}[htbp]
\centering
\includegraphics[width=0.80\linewidth]{figures/fig3_continual_retention.png}
\caption{\textbf{Phase 3 Continual Retention vs. Catastrophic Forgetting.} Full DWMA maintains $E_{A1} = 0.00500 \to E_{A2} = 0.00501$ ($\mathcal{R} = 0.9901$), while an overwriting baseline suffers catastrophic forgetting ($E_{A2} = 0.37460, \mathcal{R} = -72.97$).}
\label{fig:retention}
\end{figure}

\section{Nonlinear Structural Hypothesis Discrimination}
The environment enforces quadratic aerodynamic drag:
\begin{equation}
F_{\text{true}}(v) = -k_{\text{true}} v^2 \operatorname{sgn}(v), \quad k_{\text{true}} = 0.08
\end{equation}
The agent maintains competing linear ($H_1$) and quadratic ($H_2$) models, targeting velocities that maximize discriminative variance $\Delta(v) = |\hat{F}_1(v) - \hat{F}_2(v)|$. Model selection requires decisive log-likelihood ratio $\ln(B_{21}) > 10.0$.

\begin{figure}[htbp]
\centering
\includegraphics[width=0.80\linewidth]{figures/fig4_hypothesis_survival_analysis.png}
\caption{\textbf{Phase 4 Kaplan-Meier Survival Analysis.} Probability of remaining ambiguous $S(t) = P(T > t)$ over time. Epistemic action selection achieves median latency of $7.0$ steps (0\% censored) vs. $23.0$ steps for random babbling (28\% right-censored at $T=30$; log-rank $\chi^2 = 92.28, p < 10^{-20}$).}
\label{fig:survival}
\end{figure}

\section{Consolidated Empirical Evaluation Matrix}
Table~\ref{tab:master_matrix} compiles the complete 4-phase evaluation program across $N = 50$ seeds.

\begin{table*}[t]
\centering
\small
\caption{\textbf{Consolidated Master Empirical Evaluation Matrix ($N=50$ Seeds, $B=10,000$ Bootstrap)}.}
\label{tab:master_matrix}
\begin{tabular}{lllcccc}
\toprule
\textbf{Phase \& Phenomenon} & \textbf{Condition / Policy} & \textbf{Primary Metric} & \textbf{Mean $\pm$ SD} & \textbf{Normal 95\% CI} & \textbf{Bootstrap 95\% CI} & \textbf{Pass Rate} \\
\midrule
Phase 1: BM-1 & Zero-Shot Transfer & Gain & $85.82\% \pm 2.53\%$ & $[85.12, 86.53]\%$ & $[85.15, 86.51]\%$ & 49/50 (98\%) \\
Phase 1: BM-2 & Blind Navigation & Steps to Goal & $23.10 \pm 2.46$ st & $[22.42, 23.78]$ & $[22.44, 23.78]$ & 50/50 (100\%) \\
Phase 1: BM-3 & Silent Drift & Shock MSE & $1.384 \pm 0.120$ & $[1.351, 1.417]$ & $[1.351, 1.416]$ & 50/50 (100\%) \\
Phase 1: BM-4 & Epistemic Inquiry & Info Gain & $5.14 \pm 0.23$ nats & $[5.08, 5.21]$ & $[5.08, 5.21]$ & 50/50 (100\%) \\
Phase 1: BM-5 & Causal Repair & Post-Repair Error & $1.18 \pm 0.48$ px & $[1.05, 1.32]$ & $[1.05, 1.32]$ & 50/50 (100\%) \\
Phase 1: BM-6 & Composition & Prediction MSE & $0.119 \pm 0.075$ & $[0.098, 0.139]$ & $[0.098, 0.139]$ & 50/50 (100\%) \\
Phase 1: BM-7 & Polarity Inversion & Re-aligned $\cos(\theta)$ & $0.9915 \pm 0.0057$ & $[0.9899, 0.9931]$ & $[0.9899, 0.9931]$ & 48/50 (96\%) \\
\midrule
Phase 2: Baseline & Tracking $A \to A$ & Error $E(t)$ & $0.0050 \pm 0.0000$ & $[0.0050, 0.0050]$ & $[0.0050, 0.0050]$ & 50/50 (100\%) \\
Phase 2: Transfer & Shock $A \to B$ & Mean Shock & $1.3079 \pm 0.0421$ & $[1.2963, 1.3196]$ & $[1.2968, 1.3199]$ & Shock \\
Phase 2: Transfer & Shock $A \to C$ & Mean Shock & $0.8117 \pm 0.0272$ & $[0.8042, 0.8192]$ & $[0.8041, 0.8192]$ & Shock \\
Phase 2: Adaptation & Re-tuning $A \to B'$ & Latency $T_\epsilon$ & $16.32 \pm 0.62$ st & $[16.15, 16.49]$ & $[16.14, 16.48]$ & 50/50 (100\%) \\
Phase 2: Adaptation & Re-tuning $A \to B'$ & Rate $A_{\text{adapt}}$ & $0.0461 \pm 0.0084$ & $[0.0438, 0.0485]$ & $[0.0438, 0.0484]$ & 50/50 (100\%) \\
\midrule
Phase 3: Retention & Full DWMA & Error $E_{A1} \to E_{A2}$ & $0.00500 \to 0.00501$ & $[0.00500, 0.00502]$ & $[0.00500, 0.00502]$ & 50/50 ($R=0.99$) \\
Phase 3: Retention & Overwriter & Error $E_{A1} \to E_{A2}$ & $0.00500 \to 0.37460$ & $[0.3657, 0.3835]$ & $[0.3657, 0.3835]$ & 0/50 (Fail) \\
\midrule
Phase 4: Hypothesis & DWMA Epistemic & Log-Likelihood $\ln(B_{21})$ & $30.35 \pm 8.03$ & $[28.12, 32.57]$ & $[28.10, 32.50]$ & 50/50 (100\%) \\
Phase 4: Hypothesis & Random Babbler & Log-Likelihood $\ln(B_{21})$ & $20.02 \pm 14.34$ & $[16.05, 23.99]$ & $[16.11, 24.03]$ & 36/50 (72\%) \\
Phase 4: Survival & DWMA Epistemic & KM Median Latency & $7.0$ st & 0\% censored & $[6.58, 7.98]$ & 50/50 uncensored \\
Phase 4: Survival & Random Babbler & KM Median Latency & $23.0$ st & 28\% censored & $[19.68, 24.08]$ & Log-rank $\chi^2 = 92.28$ \\
\bottomrule
\end{tabular}
\end{table*}

\section{Conclusion}
DWMA demonstrates that rich, reusable, and adaptive predictive structure can be autonomously acquired and maintained through embodied sensorimotor interaction without relying on an LLM as its cognitive core. All four empirical phases confirm that active epistemic inference and multi-schema retention provide robust, grounded computational foundations for artificial intelligence.

\begin{thebibliography}{9}
\bibitem{friston2010} K. Friston, ``The free-energy principle: a unified brain theory?'', \textit{Nature Reviews Neuroscience}, 11(2), 127--138, 2010.
\bibitem{rao1999} R. P. Rao and D. H. Ballard, ``Predictive coding in the visual cortex: a functional interpretation of some extra-classical receptive-field effects'', \textit{Nature Neuroscience}, 2(1), 79--87, 1999.
\bibitem{ha2018} D. Ha and J. Schmidhuber, ``Recurrent world models facilitate policy evolution'', \textit{NeurIPS}, 31, 2018.
\bibitem{harnad1990} S. Harnad, ``The symbol grounding problem'', \textit{Physica D}, 42(1-3), 335--346, 1990.
\bibitem{french1999} R. M. French, ``Catastrophic forgetting in connectionist networks'', \textit{Trends in Cognitive Sciences}, 3(4), 128--135, 1999.
\bibitem{kaplan1958} E. L. Kaplan and P. Meier, ``Nonparametric estimation from incomplete observations'', \textit{JASA}, 53(282), 457--481, 1958.
\bibitem{mantel1966} N. Mantel, ``Evaluation of survival data and two new rank order statistics arising in its consideration'', \textit{Cancer Chemotherapy Reports}, 50(3), 163--170, 1966.
\end{thebibliography}

\end{document}
